\section{Performance}\label{sec:perf}

\subsection{Method and environment}
All timings come from \texttt{paper/experiments/} and are stored with their environment in \texttt{paper/results/}. Each quantity is the minimum over repeats after a warm-up call, which reduces the influence of caching and other load on the machine.
Except for the PyTorch experiment, they were measured on one x86-64 Linux machine with \envCpus{} CPU cores, Python~\envPython{}, NumPy~\envNumpy{} and SciPy~\envScipy{}, at repository commit \texttt{\envCommit{}}, without a GPU and with the default number of BLAS threads. The PyTorch experiment (\cref{sec:perf:torch}) needs a GPU and was run on a different machine, described there, so its times must not be compared with the others.
Absolute times and speed-ups depend on the machine, the BLAS library and the baseline implementation, and should be read as indicative; they can be reproduced with \texttt{run\_all.py} in \texttt{paper/experiments/}.
We compare \rom{} and the batched paths with the \emph{reference implementations in the same packages}, not with other software.

\subsection{Many-query parameter sweep}
For the two-region cantilever of \cref{sec:verif}, a sweep of \eOneSweep{} random parameter pairs takes \eOneTfom{}~ms with the full model (assembly of $\K$ and a dense solve of the free DOF for each pair) and \eOneTrom{}~ms with the affine ROM (assembly of the reduced matrix and an $r\times r$ solve), a ratio of \eOneSpeedup{}. This includes neither the offline cost of the training solves, the POD and the projections, nor the expansion back to the full field.
The ratio depends on the baseline: the full-order solves here are dense, and a banded or sparse factorisation would reduce it. The comparison shows the structure of the saving (cost independent of $n$ per query), not a universal factor.

\subsection{Batched stress recovery}
\Cref{tab:recovery} and \cref{fig:recovery} compare the batched recovery of \cref{sec:recovery} with the per-element loop for meshes of Quad4 and Hex8 elements. The speed-up lies between \recoveryMinSpeedup{} and \recoveryMaxSpeedup{} for the sizes where the loop was run, and the two paths agree to rounding error.
The loop was not run above \recoveryLoopMax{} elements because it is slow: it took between \recoveryLoopMsMin{} and \recoveryLoopMsMax{}~ms per element in these runs. The batched path processes \recoveryBigElements{} Quad4 elements in \recoveryBigMs{}~ms on this machine.

\begin{table}[t]
\centering\small
\caption{Stress recovery, batched against loop; ``--'' means the loop was not run.}\label{tab:recovery}
\begin{tabular}{lrrrrrr}
\toprule
element & elements & DOF & batched (ms) & loop (ms) & speed-up & max rel. diff. \\
\midrule
quad4 & 100 & 242 & 0.9 & 22 & 23 & $9\times10^{-16}$ \\
quad4 & 400 & 882 & 2.1 & 86 & 42 & $1\times10^{-15}$ \\
quad4 & 1600 & 3362 & 7.6 & 364 & 48 & $2\times10^{-15}$ \\
quad4 & 3600 & 7442 & 15.2 & 811 & 53 & $3\times10^{-15}$ \\
quad4 & 10000 & 20402 & 42.5 & -- & -- & -- \\
hex8 & 216 & 1029 & 4.4 & 158 & 36 & $3\times10^{-16}$ \\
hex8 & 1000 & 3993 & 19.2 & 715 & 37 & $9\times10^{-16}$ \\
hex8 & 2744 & 10125 & 48.4 & 2056 & 43 & $1\times10^{-15}$ \\
hex8 & 8000 & 27783 & 141.1 & -- & -- & -- \\
\bottomrule
\end{tabular}

\end{table}

\subsection{Hyper-reduction}
The model is a geometrically nonlinear clamped--clamped beam reduced to \ecswModes{} modes; ECSW is trained on 60 sampled states and tested on 20 others, and the trajectory is \ecswSteps{} Runge--Kutta steps. \Cref{tab:ecsw} varies the training tolerance and the number of elements.
For 120 elements and a tolerance of $10^{-4}$, ECSW keeps \ecswBigKept{} of \ecswBigTotal{} elements, evaluates the reduced force \ecswBigForceSpeed{} times faster than full assembly, with a held-out force error of \ecswBigForceErr{}, and integrates the trajectory \ecswBigTrajSpeed{} times faster with a relative trajectory error of \ecswBigTrajErr{}.
The speed-up grows with the number of elements, because the cost of the hyper-reduced force depends on the number of kept elements while the full cost grows with the mesh; at 40 elements it is smaller. The baseline is the package's own full internal-force assembly, which loops over elements in Python.

\begin{table}[t]
\centering\footnotesize
\caption{ECSW for the nonlinear beam: elements kept, held-out reduced-force error, and speed-up of one reduced-force evaluation and of a trajectory.}\label{tab:ecsw}
\begin{tabular}{rrrrrrr}
\toprule
elements & tol & kept & force err. & force speed-up & trajectory err. & trajectory speed-up \\
\midrule
40 & $1\times10^{-2}$ & 10 / 40 & $2.8\times10^{-2}$ & 5.0 & $2.0\times10^{-2}$ & 4.0 \\
40 & $1\times10^{-3}$ & 13 / 40 & $1.8\times10^{-3}$ & 3.9 & $9.8\times10^{-4}$ & 3.2 \\
40 & $1\times10^{-4}$ & 16 / 40 & $3.6\times10^{-5}$ & 3.1 & $1.1\times10^{-5}$ & 2.6 \\
40 & $1\times10^{-5}$ & 16 / 40 & $1.3\times10^{-5}$ & 3.1 & $6.7\times10^{-6}$ & 2.6 \\
120 & $1\times10^{-4}$ & 16 / 120 & $1.2\times10^{-4}$ & 9.4 & $4.5\times10^{-5}$ & 7.3 \\
\bottomrule
\end{tabular}

\end{table}

\begin{figure}[t]
\centering
\includegraphics[width=0.48\textwidth]{figures/e2_recovery_speedup.pdf}\hfill
\includegraphics[width=0.48\textwidth]{figures/e4_ecsw.pdf}
\caption{Left: recovery time against number of elements for batched and loop paths. Right: held-out reduced-force error of ECSW against the number of elements kept, for the 40-element beam (see \cref{tab:ecsw}).}\label{fig:recovery}
\end{figure}

\subsection{Component mode synthesis}
Craig--Bampton reduction shrinks the example of \cref{sec:verif} from \cmsFullSize{} to \cmsSizeFive{} DOF at a frequency error of \cmsErrFive{} (\cref{tab:cms}). We did not time this case, since the model is too small for timings to be meaningful, and make no speed claim for it.

\subsection{PyTorch backend}\label{sec:perf:torch}
\ifnum\torchAvailable=1
The optional torch paths were timed on a separate \torchOs{} machine with \torchCpus{} logical CPU cores, Python~\torchPython{}, NumPy~\torchNumpy{} and SciPy~\torchScipy{}, at commit \texttt{\torchCommit{}}. The script \texttt{e6\_torch\_backend.py} times the solve of the reduced system (\texttt{ReducedSystem.solve}) and the stress recovery on Quad4 plates of increasing size, in double precision. PyTorch \torchVersion{} was built for CUDA \torchCuda{}.
\ifnum\torchHasGpu=1 The GPU was an \torchGpu{} with NVIDIA driver \torchDriver{}.\else\relax No GPU was available on the machine that produced these results, so only CPU columns are shown.\fi{}
The conjugate-gradient solve uses a Jacobi preconditioner and a relative tolerance of $10^{-10}$; the dense solve is shown only where the system is small enough. All torch results agree with the SciPy and NumPy results to a relative difference of at most \torchMaxDiff{}.
\Cref{tab:torch} gives solve times and \cref{tab:torchstress} recovery times. On the largest system (\torchBigDof{} free DOF) the SciPy solve takes \torchBigScipyMs{}~ms and the torch CG solve on the CPU \torchBigCpuMs{}~ms, a SciPy-to-torch time ratio of \torchBigCpuRatio{}\ifnum\torchHasGpu=1\relax\space(\torchBigGpuMs{}~ms and a ratio of \torchBigGpuRatio{} on the GPU)\fi{}, where a ratio above one means torch is faster.
For stress recovery on the largest mesh, NumPy takes \torchBigStressNumpyMs{}~ms and torch on the CPU \torchBigStressCpuMs{}~ms\ifnum\torchHasGpu=1\relax\space and on the GPU \torchBigStressGpuMs{}~ms\fi{}.
Among the sizes tested, torch CG on the CPU first beats SciPy at \torchCrossCpu{} free DOF and on the GPU at \torchCrossGpu{} (``none'' means never), and torch stress recovery first beats NumPy at \torchStressCrossCpuElems{} elements on the CPU and \torchStressCrossGpuElems{} on the GPU; at the smaller sizes tested the torch path was slower. The crossover points are only resolved to the coarse grid of sizes in the tables.
The GPU is a consumer card with limited double-precision throughput and all computations are in double precision, so results on other GPUs will differ. The purpose of the torch backend is differentiability and device placement rather than speed; these numbers show what it costs and where it starts to pay off, and they depend on the PyTorch build and hardware.

\begin{table}[t]
\centering\footnotesize
\caption{Solve times in ms for the reduced system of a Quad4 plate: SciPy against torch (CG and dense) on each available device; ``--'' means not run.}\label{tab:torch}
\begin{tabular}{rrrrrr}
\toprule
free DOF & SciPy & torch CG cpu & torch dense cpu & torch CG cuda & torch dense cuda \\
\midrule
1088 & 4.8 & 25.6 & 11.6 & 66.8 & 40.4 \\
4224 & 47.2 & 92.9 & 436.4 & 234.3 & 879.5 \\
16640 & 363.1 & 247.1 & -- & 383.4 & -- \\
66048 & 2070.4 & 2284.7 & -- & 836.0 & -- \\
\bottomrule
\end{tabular}

\end{table}

\begin{table}[t]
\centering\footnotesize
\caption{Stress-recovery times in ms: batched NumPy against torch on each available device.}\label{tab:torchstress}
\begin{tabular}{rrrr}
\toprule
elements & NumPy & torch cpu & torch cuda \\
\midrule
512 & 3.9 & 5.5 & 17.4 \\
2048 & 10.0 & 10.6 & 17.0 \\
8192 & 58.1 & 27.1 & 18.2 \\
32768 & 165.7 & 70.8 & 47.7 \\
\bottomrule
\end{tabular}

\end{table}
\else
The torch solve and the differentiable recovery are covered by tests (\cref{sec:arch}), but the timing experiment \texttt{e6\_torch\_backend.py} could not run on the machine above because PyTorch is not installed there. \todo{run e6 on the CPU and GPU machine, then \texttt{make\_tables.py}; this paragraph, a table and a short discussion are generated from the stored result.}
\fi

\subsection{Limits of these measurements}
Two machines (one CPU-only, one with a consumer GPU for the torch experiment), small and medium problems, no comparison with other packages, and baselines that are the packages' own reference paths. The ROM ratio excludes offline cost. Large-scale performance (millions of DOF, parallel runs) is outside what the scripts measure and outside the aims of the package.
